\documentclass[11pt]{letter}
\usepackage[margin=1in]{geometry}
\usepackage[T1]{fontenc}
\usepackage{lmodern}
\usepackage{hyperref}

\signature{Meng Li, on behalf of Xiaohua Yang, Jie Liu, and Shiyu Yan\\
School of Computing, University of South China\\
Hengyang 421001, China\\
\href{mailto:mlemon@usc.edu.cn}{mlemon@usc.edu.cn}}

\begin{document}

\begin{letter}{Professor Mauro Pezze\\
Editor-in-Chief\\
ACM Transactions on Software Engineering and Methodology}

\opening{Dear Professor Pezze,}

We are pleased to submit our manuscript ``NOETHER: Constructive Metamorphic Pattern Identification from Operator Algebras and a Falsifiable Invariance-Blindness Theorem'' for consideration in ACM TOSEM under the \textbf{Fast Impact} track.

The paper targets a software-testing problem that directly concerns the TOSEM readership: metamorphic testing is valuable when test oracles are unavailable, but metamorphic-relation (MR) identification remains largely experience- or search-driven. Existing approaches offer limited structural explanation of where MR patterns come from, when a pattern catalogue is closed, or how it transfers across program families. NOETHER addresses this origin--closure--transferability gap by turning MR identification into an auditable construction over a program family's operator-algebraic structure.

The manuscript makes five contributions to software-engineering methods and foundations:

\begin{enumerate}
\item A layered, constructive framework (NOETHER) that derives MetaPatterns from the operator-algebraic structure of a program family by the mechanical CONSTRUCT-MP algorithm.
\item A no-drop closure invariant and polynomial-time constructibility result over the algebra-induced MR space.
\item A falsifiable Invariance-Blindness Theorem characterising exactly which faults symmetry and self-adjoint MetaPattern MRs can and cannot detect within the stated fault class.
\item A negative PWR reactor-physics instantiation showing that absolute completeness is false and identifying five independent obstructions.
\item A structural-transferability and evidence protocol instantiated on Boltzmann reactor physics, equivariant machine learning, and relational query optimisers.
\end{enumerate}

The evidence combines analytical and experimental validation paradigms. The closure and constructibility claims are supported by formal derivations; the blindness theorem is stated as a falsifiable prediction about detectable and undetectable fault classes; the PWR instantiation deliberately reports a negative result to delimit the theory; and the three-domain protocol checks whether the construction transfers beyond a single case. The submitted data and artifact availability statement identifies the replication package, supplementary material (S1--S12), prompts, raw outputs, mapping files, and scripts needed to inspect or rerun the reported checks.

The manuscript has been shortened so that the main paper body fits the TOSEM Fast Impact 45-page limit in the official ACM \texttt{acmsmall,screen,review} format. References and online-only appendix/supplementary material are kept outside the main-body count. Self-contained implementation details, detailed derivation traces, and secondary tables have been migrated to appendices and supplementary material.

A preprint of this work is available on arXiv (arXiv:2605.17390, \url{https://arxiv.org/abs/2605.17390}). TOSEM is a single-blind journal, so the submitted manuscript lists all authors and affiliations on the first page; this cover letter discloses the related preprint for editorial transparency.

A separate line of work by a subset of the authors addresses the orthogonal problem of selecting a minimum complete subset from a given MR pool. That work takes an MR set as input and minimises its cardinality under a fixed fault model. The present paper concerns the upstream question of where MRs and MetaPatterns come from: their constructive derivation from a program-induced operator algebra and algebraic closure under the \texttt{Translate} operator. The two share no theorem or empirical claim.

In line with ACM's Policy on Authorship, large language models appear in this work in two roles. As instruments of the study, they are part of the reported LLM baseline, LRCA second-rater protocol, and mutant-equivalence adjudication; the prompts and raw outputs are released with the replication package. As authoring assistance, the authors used a large-language-model assistant for code scaffolding, reference formatting, and language editing. All research design, theorems, proofs, experimental procedures, numerical results, and claims were produced and verified by the authors, who take full responsibility for the manuscript.

The authors declare no competing interests.

\closing{Sincerely,}

\end{letter}

\end{document}
